\chapter{Proofs for derivable priorities and applicability}
\label{annex:p18-is-dynamic}

\begin{annexitem}{annex:p18.is-dynamic-stability}{thm:p18-is-dynamic-stability}{Compatible support and stability}
\begin{proof}
Suppose \(E_X\) contained both \((r,s)\) and \((s,r)\).
An argument in \(X\) concluding the first has a strict extension
concluding \(\neg(s\triangleleft r)\), which conflicts with the
argument concluding the second. For \(r=s\) the same reasoning
gives self-conflict. Both possibilities contradict conflict-freedom.

If \(R<_XR'\) and \(R'<_XR\), choose their existential witnesses
\(r\in R\), \(s\in R'\). They give both \((r,s)\) and \((s,r)\),
impossible. If either set is empty, one of the two existential
clauses already fails. This proves lift asymmetry without
transitivity. For a complementary pair of strict-extension
conclusions, the two frontier comparisons therefore cannot both
hold. At least one rebuttal direction survives. If that direction
is blocked by a reverse undercut, the reverse undercut itself is
a defeat. An undercut is already a defeat. These observations prove
orientation and the attack qualification.

For \(X\subseteq Y\), every positive frontier comparison true at
\(X\) remains true at \(Y\). Rebuttal can therefore disappear but
cannot appear; the undercut and incoherence clauses do not change.
Thus \(D_Y\subseteq D_X\).
Suppose \(A D_X^s B\). If its forward defeat is an undercut or
the empty incoherence clause, it persists at \(Y\); the reverse
defeat cannot appear, by the just-proved inclusion.
Otherwise it has a rebuttal witness, and neither argument undercuts
the other. In particular, the absence of \(B D_XA\) says that
the reverse rebuttal fails for every complementary witness.
For the chosen witness with frontiers \(R_A,R_B\), this means
\(R_B<_XR_A\). It remains true at \(Y\), so lift asymmetry implies
\(R_A\not<_YR_B\). The original forward rebuttal survives.
Its reverse still cannot appear. This proves strict stability.
\end{proof}
\end{annexitem}

\begin{annexitem}{annex:p18.is-dynamic-fixed-point}{thm:p18-is-dynamic-fixed-point}{The compatible-support fixed point}
\begin{proof}
Let conflict-free \(X\subseteq Y\), and let \(A\in G(X)\).
Every \(B D_YA\) also satisfies \(B D_XA\), hence has a
strict defender \(C\in X\). Strict stability makes the same \(C\)
a strict defender at \(Y\). Thus \(G(X)\subseteq G(Y)\).

Suppose \(A,A'\in G(X)\) conflict. By orientation, say \(A D_XA'\).
The acceptability of \(A'\) gives \(B\in X\) with \(B D_X^s A\).
The acceptability of \(A\) then gives \(B'\in X\) with
\(B'D_X^s B\). This is an attack between members of \(X\),
unless it is the empty incoherence defeat. That exception would
make \(B\) incoherent, also impossible for a member of a
conflict-free set. Thus \(G(X)\) is conflict-free.

The iteration starts increasing since \(J_0=\varnothing\).
At successors monotonicity preserves inclusion. A union at a
limit is conflict-free because any two members already occur
together at some earlier stage of the chain. For an element of
that union, its earlier acceptability and monotonicity imply its
membership in the image of the union. This proves the next
inclusion and permits transfinite continuation.
There cannot be more strict increases than distinct arguments:
choose a newly entering argument at each strict increase;
this would otherwise inject an ordinal of cardinality greater
than \(\mathcal A\) into \(\mathcal A\).
At the first equality the value is a fixed point.
If \(K\) is any conflict-free fixed point, transfinite induction
using monotonicity and unions gives \(J_\alpha\subseteq K\).
Hence the stabilized value is least.
If a fixed point of the unrestricted formula were conflict-free
but did not contain \(J\), it would contradict this leastness.
Every subset of conflict-free \(J\) is conflict-free, so no proper
subset of \(J\) can be a fixed point either. These statements use
fixed points, not an additional maximality condition.

For subargument closure, let \(A\in J\), \(A'\subseteq A\),
and \(B D_JA'\). If this is the empty incoherence case, the
inherited inconsistency makes \(A\) incoherent too, so \(B D_JA\)
and the acceptability of \(A\) supplies a strict defender.
Otherwise \(B\) attacks \(A'\), and thus attacks \(A\).
If \(B D_JA\), use the same acceptability. If not, orientation
gives \(A D_JB\), now strictly, so \(A\) itself is a defender.
Every defeater of \(A'\) has been strictly defended by \(J\).
Therefore \(A'\in G(J)=J\).

If strict arguments are inconsistent, every argument is incoherent
by strict-extension closure. The empty argument defeats every
argument, including itself. Thus none is acceptable to the empty
set. The iteration remains empty. This also explains why the
designated empty defeat must be retained in the construction.
\end{proof}
\end{annexitem}

\begin{annexitem}{annex:p18.is-dynamic-finitary}{thm:p18-is-dynamic-finitary}{Finite priority information}
\begin{proof}
Write \(\mathcal D(A)\) for the finite set of defeasible rules in
\(A\). A strict extension adds no defeasible rule, so every frontier
in an extension of \(A\) is a subset of \(\mathcal D(A)\).
There are finitely many such subsets, including the empty one.
Thus every priority-sensitive part of a defeat between \(A,B\)
depends only on membership in \(E_X\) of pairs in
\[
 [\mathcal D(A)\times\mathcal D(B)]
 \ \cup\ [\mathcal D(B)\times\mathcal D(A)].
\]
The existence of a strict extension realizing a frontier pair is
independent of \(X\); even infinitely many extensions introduce no
further comparison bits.

For an increasing sequence \(X_n\) with union \(X\), every one
of these finitely many bits present in \(E_X\) is witnessed by
an argument in some \(X_n\). Choose the maximum of those finitely
many indices. Absent bits are absent at every stage. Hence both
directions of defeat and strict defeat stabilize to their values
at \(X\).

Put \(U=\bigcup_nJ_n\). The preceding theorem gives \(U\subseteq G(U)\).
Let \(A\in G(U)\). Every potential defeater belongs to the finite
set
\(\mathcal B=\{B:B D_\varnothing A\}\):
defeat only implies attack with the possible extra empty argument,
and the theory is finitary. Each \(B\in\mathcal B\) either is no
longer a \(D_U\)-defeater or has some \(C_B\in U\) with
\(C_BD_U^sB\). In the first case finite-bit stabilization eventually
removes it; in the second choose a stage containing \(C_B\) and
large enough to stabilize that strict comparison. There are only
finitely many \(B\), so one finite \(N\) works for all cases.
It follows that \(A\in G(J_N)=J_{N+1}\subseteq U\).
Thus \(U\) is a conflict-free fixed point. Leastness and the
containment of all iterates in \(J\) give \(U=J\).
\end{proof}
\end{annexitem}

\begin{annexitem}{annex:p18.is-dynamic-game-sound}{thm:p18-is-dynamic-game-sound}{Soundness of the single-argument game}
\begin{proof}
Induct upward on the finite tree. A proponent leaf has no
\(D_\varnothing\)-defeater and therefore belongs to \(J_1\).
At a proponent node \(A\), consider any \(B D_JA\).
The inclusion \(D_J\subseteq D_\varnothing\) makes \(B\) one
of the opponent children covered by the tree. Its chosen response
\(C\) belongs to \(J\) by induction. In particular \(\{C\}\)
is conflict-free.
If \(C D_{\{C\}}^sB\), strict stability transfers this to \(J\).
If the response instead neutralizes \(B\), then
\(\neg(B D_{\{C\}} A)\), and antitonicity implies
\(\neg(B D_JA)\), contradicting the selected challenge.
Thus every actual \(J\)-defeater has a strict defender in \(J\);
\(A\in G(J)=J\). The proof does not need a completeness assumption,
minimal arguments, or a merger of responses.
\end{proof}
\end{annexitem}

\begin{annexitem}{annex:p18.is-dynamic-certificates}{thm:p18-is-dynamic-certificates}{Support-explicit correspondence and finite computation}
\begin{proof}
For soundness induct on certificate height.
All arguments in a support \(X_B\) belong to \(J\) by their
strictly smaller certificates. Neutralization at \(X_B\)
therefore remains neutralization at \(J\); strict defense at
\(X_B\) remains strict defense at \(J\).
An actual \(D_J\)-attacker cannot take the first alternative.
Every such attacker is a \(D_\varnothing\)-attacker and is
covered, so the claimed argument belongs to \(G(J)\).
The leaf case is immediate from \(G(\varnothing)\).

For completeness in a finitary theory, use induction on the
least positive integer \(n\) for which \(A\in J_n\).
At \(n=1\) there are no \(D_\varnothing\)-attackers.
For \(n>1\), enumerate the finite set of such attackers \(B\).
If \(B\) no longer defeats \(A\) relative to \(J_{n-1}\),
the finite-information proof selects finitely many arguments
from \(J_{n-1}\) witnessing all present comparison bits needed
to give that same outcome. Their set is a valid \(X_B\).
If \(B\) still defeats \(A\), select its strict defender
\(C_B\in J_{n-1}\), and finitely many witnesses for the bits
of both comparisons between \(C_B\) and \(B\).
Include \(C_B\) in \(X_B\). All other present bits introduced
by these arguments also occur in \(E_{J_{n-1}}\); absent bits
cannot be introduced. The finite comparison outcome is therefore
unchanged. Each \(X_B\) is a subset of conflict-free \(J_{n-1}\).

Every member of every \(X_B\) has smaller entry rank than \(A\)
and has a certificate by induction. There are finitely many
attackers and finite supports, so their attached certificates
give a finite construction. Strict descent of the entry rank
also prevents repetition of a claim on a descending branch.
This proves completeness without combining support arguments.

For finitely many ground rules, legal content is a subset of a
finite set, and legality is decided by checking the existence
of a permitted ordering with unique heads. Strict saturation
and frontiers are computed on finite rule data as in
Theorem~\ref{thm:p18-is1996-finite}. Each \(G\)-evaluation is
therefore a finite table calculation. Its increasing iteration
has at most \(|\mathcal A|\) strict increases.
The preceding inductive proof extracts finite certificates
from that table. No effective infinite-domain claim is used.
\end{proof}
\end{annexitem}

\begin{annexitem}{annex:p18.is-dynamic-weak-game}{thm:p18-is-dynamic-weak-game}{Reversed roles on a finite argument universe}
\begin{proof}
Fix a finite universe and write \(D,J_n,J\) for the fixed-priority
objects. Suppose some member of \(J\) defeats \(A\). Select such
a member \(C\) of least entry rank \(n\). It strictly defeats
\(A\): otherwise \(A D C\), and acceptability of \(C\) at
stage \(n\) gives \(C'\in J_{n-1}\) strictly defeating \(A\),
contrary to minimality of \(n\).
For \(n=1\) the same conclusion follows because \(C\) has no
defeater at all.

The reversed-game opponent can open with this strict defeater
\(C\), then use the strong proof construction of
Theorem~\ref{thm:p18-is1996-proof-correctness} with roles
interchanged. Its own moves have strictly decreasing entry
ranks, hence cannot repeat, and its replies are strict defeats.
The other player's moves are ordinary defeats, exactly the
challenges covered by that construction. It thus wins.

Conversely, if the opponent has a winning strategy, its first
move \(C\) strictly defeats \(A\). The subtree after \(C\),
with players interchanged, is a finite strong proof tree for
\(C\): it covers all ordinary defeating responses, and its
proponent replies are strict with no repetition. Soundness
puts \(C\) in \(J\).
Every play of the reversed game is finite, since each opponent
move consumes a new element of a finite universe and moves
alternate. Finite branching and the uniform bound of
\(|\mathcal A|\) opponent moves permit backward induction.
At a position with no legal move the player to move loses;
otherwise a player wins exactly when one successor is winning
for that player. Thus precisely one player has a winning strategy,
proving the equivalence.

If a justified \(C\) attacks \(A\), orientation gives \(CDA\)
or \(ADC\). In the second case acceptability of \(C\) supplies
a justified strict defeater of \(A\). Hence attack-overruling
implies defeat by \(J\), giving the stated one-way consequence.
For the counterexample, the legal contents are
\([], [r], [s], [r,s]\). The singletons defeat each other and
the last content is incoherent. Only the empty argument is
unattacked at the first stage. It does not defeat either
singleton, so subsequent stages add neither. It does defeat
the incoherent combined content, which therefore never enters.
The empty argument has no conclusion or assumption and attacks
nothing, yet its designated defeat of the combined content is
strict. This proves every assertion in the residual example.
\end{proof}
\end{annexitem}

\begin{annexitem}{annex:p18.is-dynamic-examples}{thm:p18-is-dynamic-examples}{Priority examples against all generated challenges}
\begin{proof}
First consider the outside-domain example. Every attacker of
\([r_1]\) must use \(r_3\), hence \(r_2\).
At \(X\), \([r_4]\) strictly defeats all those contents through
the \(b,\neg b\) conflict, because \(r_2\triangleleft r_4\)
is asserted and its reverse is not. This proves membership
in \(G(X)\). At \(Y\), take the particular attacker
\(B=[r_2,r_3]\). It still defeats \([r_1]\), since no priority
relates \(r_3\) to \(r_1\).
Neither priority singleton in \(Y\) attacks \(B\).
The remaining possible defender \([r_4]\) no longer rebuts \(B\):
its frontier is now also lower than \(\{r_2\}\).
It does not undercut \(B\). Thus \(B\) has no strict defender
in \(Y\), proving the asserted failure of monotonicity.

In the reinstatement example the only priorities have endpoints
in the two separate components \(\{r_0,r_3\}\) and
\(\{r_4,r_5\}\). The default \(r_6\) is unattacked.
An attack on \([r_4]\) must contain \(r_5\), whose frontier
is below \(\{r_4\}\) once \(r_6\) has entered.
Thus \([r_4]\in J_2\).
Every attack on \([r_3]\) requires \(r_0\), now lower than
\(\{r_3\}\), so \([r_3]\in J_3\).
Every attack on \([r_2]\) is an undercut producing \(b\).
There is no producer of \(\neg c\).
The unique possible producer of \(b\) is \(r_1\), with
antecedent supplied by \(r_0\). The accepted \([r_3]\)
strictly defeats every such content; reverse rebuttal of
\([r_3]\) must use the lower \(r_0\).
Hence \([r_2]\in J_4\). Additional priority conclusions in
an attacker do not change the comparison support \(J_3\).

For the three-cycle example, \(r_7,r_8\) concern a component
of names disjoint from \(\{r_1,r_2,r_3\}\).
No reverse comparison in that component can be produced.
Their separate and combined arguments are unattacked.
Let \(C=[r_4,r_5]\). Its strict closure is the acyclic chain
\(r_1\triangleleft r_2\triangleleft r_3\), its transitive edge,
and the reverse negative edges. No argument using only
\(r_4,r_5,r_7,r_8\) conflicts with this chain: the strict
rules preserve these two acyclic components.
Thus every attack on \(C\) has a witnessing frontier containing
\(r_6\). At \(J_1\), this witness is lower than every nonempty
frontier of \(C\), since those frontiers are subsets of
\(\{r_4,r_5\}\) and both required comparisons are asserted.
There is no empty strict-fact frontier in this example.
Every reverse rebuttal is therefore blocked. There are no
undercuts and \(C\) is coherent. Hence \(C\) has no
\(D_{J_1}\)-defeater and belongs to \(J_2\).
This argument includes the attackers combining \(r_6\)
with either or both lower-chain defaults; it does not rely
on a graph containing only five singleton arguments.

In the five-rule neutralization example, the only source of a positive priority is
\(r_5\); its strict consequences contain its reverse negative,
but not its own negation. No rule concludes a contrary to any
conclusion of \([r_5]\), and there are no weak antecedents.
Thus \([r_5]\) is unattacked and belongs to \(J_1\).
Any coherent attack on \([r_4]\) requires \(b\), whose unique
possible producer is \(r_2\). At support \(J_1\) its frontier
is lower than \(\{r_4\}\); it cannot defeat \([r_4]\).
An incoherent attacker cannot create a surviving reverse
undercut, and the same frontier comparison blocks its \(b\)
rebuttal. There is no empty incoherence defeat of \([r_4]\).
Hence \([r_4]\in J_2\).
Every attacker of \([r_1]\) must produce \(\neg a\) using
\(r_3\), and must therefore contain \(r_2\).
The argument \([r_4]\) strictly defeats each such content:
its \(\neg b\) rebuttal survives, whereas any reverse \(b\)
rebuttal uses the lower \(r_2\). There are no other conclusions
in \([r_4]\) to attack. Thus \([r_1]\in J_3\).
These observations also justify the recorded self-context
defense and neutralization moves.

For the empty-antecedent variant of the circular example,
the priority default is unattacked; at its established comparison
every attack on \([r_1]\) requires the lower \(r_2\).
Thus the priority enters first and \([r_1]\) at the next stage.
For the conditional circular example every nonempty argument contains
\(r_1\) or \(r_2\). If it contains both it is incoherent and is
defeated by the empty argument. Otherwise the other singleton
defeats it at empty comparison: neither has weak premises and
there is no established priority. Therefore \(J_1=\{[]\}\).
The coherent singleton attackers of these contents are not
defeated by the empty argument; thus no further content enters.
The asserted fixed point follows. For the two opposed priority
defaults, the asymmetry rule lets each singleton strictly extend
to the negation of the other's conclusion. At empty comparison
each defeats the other; their combination is incoherent.
The same argument gives \(J=\{[]\}\).

In the building example the five ordinary factual arguments and
the higher-order factual argument have no contrary producers.
The only positive comparisons generated by defeasible priority
rules are
\[
R_M\triangleleft R_P,\qquad R_P\triangleleft R_M,\qquad
T_1\triangleleft K_1.
\]
The last comparison is unattacked: the other comparisons have
endpoints \(R_M,R_P\), disjoint from \(T_1,K_1\), and the strict
priority rules cannot connect these two components.
Their negative conclusions are precisely the ones forced within
those components. An argument for \(T_1\triangleleft K_1\),
including its factual premise, is therefore in \(J_1\).

Every attack on the argument using \(K_1\) must use the reverse
lower comparison and therefore \(T_1\); no factual premise has
a contrary producer. Its responsible frontier contains \(T_1\)
for the reverse comparison and is lower than \(\{K_1\}\)
once the higher priority is established.
Arguments combining both lower comparisons are incoherent.
The argument using \(K_1\) is consequently accepted at the next
stage. Every attack on the preservation argument then requires
\(R_M\) to conclude modification. The established lower priority
suppresses that defeat. This yields preservation at a further
stage. In the displayed dialogue the \(K_1\) argument neutralizes
the initial modification challenge, the \(T_1\) argument challenges
its comparison under empty support, and the \(T_2\) argument
neutralizes that challenge. Superarguments cannot introduce a
new contrary fact or a new positive comparison; the preceding
producer analysis covers all of them.

In the exception example, before adding \(h\), all factual
arguments and \(q\) are unattacked. The only contrary producer
for the sale conclusion is \(L_0\); its responsible frontier
is lower than \(\{E_0\}\) after \(q\) enters.
No argument produces either of the two exception atoms.
Thus the \(E_0\) argument is accepted and attacks every
argument using \(L_0\), which is overruled.
After adding \(h\), its singleton is unattacked and undercuts
every argument using \(E_0\). There is no reverse undercut of
\([h]\), and no contrary to its conclusion, so this defeat is
strict independently of all priorities.
Every challenge to the \(L_0\) argument must use \(E_0\);
\([h]\) strictly defeats each such challenger. The local
argument is accepted at the next stage. Its contrary sale
argument remains overruled through the justified exception.
The finite ground-instance restriction in the body prevents
unmentioned instances of the displayed schemata from adding
different conclusions.
\end{proof}
\end{annexitem}

\begin{annexitem}{annex:p18.is-dynamic-applicability}{thm:p18-is-dynamic-applicability}{Applicability constructions and countermodels}
\begin{proof}
In the negative translation, the translated rule has the weak
antecedent \(\mathord{\sim}\neg\operatorname{app}(r)\).
An argument concluding \(\neg\operatorname{app}(r)\) therefore
undercuts each content using it, by the definition of \(U\).
In the positive translation an argument must first produce
\(\operatorname{app}(r)\); its default producer is \(\rho_r\).
At a conflict on that conclusion the responsible frontier is
\(\{\rho_r\}\), not \(\{r\}\).
Both translations use only the existing rule language and
defeat definition. The general core theorems apply directly,
without any assertion that the two translations are equivalent.

For priority-free direct exclusion, an exclusion attack from
\(A\) to \(B\) is a defeat unless \(B\) undercuts \(A\);
in that case \(B\) defeats \(A\). Thus orientation holds for
the enlarged attack relation. Such an attack on a subargument
persists against its superargument, since the excluded rule
still occurs. Defeat still implies enlarged attack except for
the empty case.
The enlarged conflict-free condition implies the original one,
so extraction of priorities remains asymmetric.
Exclusion is independent of \(X\), and adding support can only
remove rebuttals. A strict defeat given by exclusion persists;
no reverse defeat can appear. For a strict defeat supplied only
by rebuttal, the absence of reverse defeat again implies absence
of reverse rebuttal, so the same frontier proof gives stability.
These are exactly the premises used in the fixed-point proof.
For finite-information and certificate proofs, exclusion adds
no comparison bit. The same arguments apply with enlarged
attackers in the finitary hypothesis. Strict inconsistency
still propagates by strict extensions; exclusion targets no
strict rule. This establishes the asserted inherited results.

For the first countermodel, in the direct version the legal
contents are \([], [r], [s], [r,s]\).
The argument \([r]\) undercuts every content using \(s\).
Therefore exclusion by such a content cannot defeat \([r]\),
and the combined content is incoherent. No argument defeats
\([r]\), so it is justified at the first stage.
In the negative-premise version the two singletons undercut
each other, and their combined content is again incoherent.
Only the empty argument enters initially, and it defeats
neither coherent singleton. The iteration adds neither.
Thus the conclusions differ as claimed.

For the sensitive-exclusion example, \([s]\) and \([k]\)
do not attack each other or themselves. Thus
\(X=\{[s],[k]\}\) belongs to the specified domain.
Any exclusion attack on \([r]\) requires the only producer
\(s\) of \(\neg\operatorname{app}(r)\); its frontier is
\(\{s\}\), lower than \(\{r\}\) at \(X\).
There is no producer of \(\neg p\), no weak antecedent and
no incoherence in \([r]\). Hence \([r]\) has no
\(D_X^{\rm ex}\)-defeater. Likewise there is no contrary or
exclusion producer against \([s]\), so it too has none.
Both belong to \(G(X)\), but \([s]\) exclusion-attacks
\([r]\) by the definition that is independent of priority.
This proves failure of the claimed domain preservation,
not just failure of a particular proof.
Indeed \([s]\) and \([k]\) also enter at the empty-support
stage, so the obstruction is relevant to iteration from below.

In the final positive-premise alternative, the argument using
\(k\) has no contrary producer and enters first.
The default \(\rho_r\) then rebuts \(s\) with the established
preference in its favor. No other producer attacks its conclusion.
The rule \(r\) can use that accepted applicability support to
derive \(p\); its attacks all originate in \(s\) and are
strictly defeated by the applicability default.
Thus the positive translation realizes the stated ranked
applicability outcome while overruling \(s\).
The failing direct version accepts \(s\), so an equivalence
between those systems would be false.
\end{proof}
\end{annexitem}
